\documentclass[11pt]{article}
\usepackage[utf8]{inputenc}
\usepackage[spanish,es-nodecimaldot,es-noshorthands]{babel}
\usepackage{amsmath,amssymb,amsthm}
\usepackage{booktabs}
\usepackage{graphicx}
\usepackage[margin=2.6cm]{geometry}
\usepackage{natbib}
\usepackage[colorlinks=true, citecolor=blue, linkcolor=blue, urlcolor=blue]{hyperref}

\newtheorem{assumption}{Supuesto}
\newcommand{\E}{\mathbb{E}}
\newcommand{\Prob}{\mathbb{P}}
\newcommand{\ind}[1]{\mathbf{1}\{#1\}}

\title{Un modelo estocástico de las cotizaciones previsionales chilenas\\
\large Nota metodológica — borrador para discusión}
\author{Carlos Pérez G.\thanks{Universidad Andrés Bello. Borrador; no citar.}}
\date{\today}

\begin{document}
\maketitle

\section{Objetivo y hechos estilizados}

Esta nota especifica un proceso estocástico para la historia de cotizaciones
obligatorias de un afiliado al sistema de pensiones chileno, estimable con la
muestra de Historias Previsionales de Afiliados (HPA) de la Superintendencia
de Pensiones. El proceso está diseñado para servir como insumo exógeno de un
modelo de ciclo de vida de inversión previsional, por lo que se privilegia una
representación markoviana explícita.

La especificación responde a cuatro hechos documentados en la muestra HPA
(panel persona--mes 1981--2023, $N=28.524$ individuos, 7,5 millones de
observaciones):

\begin{enumerate}
\item \textbf{Densidad intermedia y dispersa.} La densidad de cotización
agregada es 0,535. La distribución de densidades individuales de vida es
fuertemente bimodal (en forma de U): 15,5\% de los afiliados presenta densidad
inferior a 0,10 y 16,5\% superior a 0,90.
\item \textbf{Heterogeneidad de niveles.} Clasificando por terciles de
densidad, la probabilidad mensual de reingreso desde una laguna varía en un
factor de 8 entre grupos (0,023 vs.\ 0,178) y la persistencia mensual en
cotización va de 0,837 a 0,977.
\item \textbf{Dependencia de duración genuina y de muy largo alcance.} Los
hazards de salida de ambos estados decrecen fuertemente con la duración del
episodio \emph{dentro de cada grupo de adhesión} (razón entre el hazard a
1--3 meses y a 19--24 meses de entre 4 y 8), y la caída \emph{continúa} en
duraciones largas: el hazard de reingreso cae en un factor adicional de
$\sim$5 entre los 24 meses y los 120+ meses de laguna, a toda edad. Las
lagunas muy largas no son marginales: a los 50--65 años, la celda de
lagunas de más de 120 meses es la más poblada del estado laguna. La
dependencia de duración no es, por tanto, un artefacto de composición,
aunque la heterogeneidad también es de primer orden.
\item \textbf{Ausencia de \emph{scarring} salarial.} El salario real de
reingreso tras una laguna no muestra penalización respecto del salario
pre-laguna, una vez ajustado por el crecimiento salarial contrafactual de los
cotizantes continuos. La laguna destruye cotizaciones, no salario observado
de reentrada.
\end{enumerate}

Los hechos 2 y 3 excluyen tanto un proceso de Markov homogéneo de dos estados
(sin memoria de duración) como una explicación puramente composicional. El
modelo combina, en consecuencia, \emph{tipos discretos} y
\emph{semi-Markov con dependencia de duración truncada}.

\section{El proceso de participación}

\subsection{Estados, tipos y notación}

El tiempo es discreto y mensual, $t = 1, 2, \dots$; el individuo $i$ entra al
sistema en el mes de su afiliación y se caracteriza por su sexo $g_i$ y un
\emph{tipo} permanente $k_i \in \{1,\dots,K\}$ con probabilidades
poblacionales $\pi_{k|g}$ \emph{condicionales al sexo}. Los tipos se
denominan \emph{estable} (tipo I), \emph{intermitente} (tipo II) y
\emph{tipo III}, en orden decreciente de densidad posterior media. Los dos
primeros llevan nombre porque su firma dinámica admite un resumen de una
frase; el tercero se deja deliberadamente sin nombre: su firma es bimodal
en ambos estados —combina las lagunas más persistentes de la población con
una minoría de empleos muy largos que concentra la mayor parte de sus meses
cotizados— y cualquier etiqueta de nivel (``frágil'', ``de baja
densidad'') induciría a leer un resultado donde lo que hay es una dinámica
(Sección~\ref{sec:firmas}). Las etiquetas describen patrones de transición,
no tramos de densidad. Todo el modelo se
especifica y estima separadamente por sexo: las brechas observadas (densidad
mediana de vida de 0,644 en hombres vs.\ 0,457 en mujeres; brecha de densidad
por edad creciente de 12 a 20 puntos porcentuales entre los 25 y los 64
años) y las asimetrías institucionales (edades legales de retiro, tablas de
mortalidad) hacen inadecuada una especificación unisex con desplazamientos
aditivos. El estado de participación es
\[
s_{it} \in \{0,1\}, \qquad 1 = \text{cotiza}, \quad 0 = \text{laguna},
\]
y la duración en el estado corriente, truncada en $\bar D$ meses, es
\[
d_{it} = \min\bigl\{\#\{\tau \le t : s_{i\tau'} = s_{it} \ \forall \tau' \in
[\tau, t]\},\ \bar D\bigr\}.
\]
La duración entra en los hazards a través de una función escalonada sobre
los tramos $\mathcal{D} = \{1, \dots, 12, 13\text{--}18, 19\text{--}23,
24\text{--}35, 36\text{--}59, 60\text{--}119, 120+\}$, constante dentro de
cada tramo y desde los 120 meses ($\bar D = 240$ como tope operativo de la
simulación). Los tramos finos hasta el año y gruesos después equilibran la
resolución donde la caída del hazard es rápida (hecho 3) con el tamaño de
las celdas; una versión anterior truncaba en $\bar D = 24$ bajo el supuesto
de hazard plano desde los dos años, supuesto que los datos rechazan (el
hazard sigue cayendo un factor de $\sim$5 hasta los 120+ meses) y cuya
corrección resultó cuantitativamente importante para las cotizaciones
esperadas a edades altas.

\subsection{Hazards de transición}

Las transiciones entre estados están gobernadas por hazards en tiempo
discreto que dependen del tipo, la duración y la edad $a_{it}$, con todos los
parámetros específicos al sexo:
\begin{align}
\lambda^{01}_{k,g}(d, a) &= \Prob\left(s_{i,t+1} = 1 \mid s_{it} = 0,\
d_{it} = d,\ a_{it} = a,\ k_i = k,\ g_i = g\right), \\
\lambda^{10}_{k,g}(d, a) &= \Prob\left(s_{i,t+1} = 0 \mid s_{it} = 1,\
d_{it} = d,\ a_{it} = a,\ k_i = k,\ g_i = g\right).
\end{align}
Se especifican como logits con perfiles de duración específicos al
macro-tramo de edad $B(a) \in \{18\text{--}34,\ 35\text{--}49,\
50\text{--}65\}$:
\begin{equation}
\lambda^{sr}_{k,g}(d, a) = \Lambda\bigl(\alpha^{sr}_{k,g,B(a)} +
f^{sr}_{k,g,B(a)}(d) + h^{sr}_{k,g}(a)\bigr), \qquad
\Lambda(x) = \frac{e^x}{1+e^x},
\label{eq:logit}
\end{equation}
donde $\alpha^{sr}_{k,g,B}$ es el intercepto específico al tipo, sexo,
estado de origen y macro-tramo, $f^{sr}_{k,g,B}(\cdot)$ es la función
escalonada de la duración sobre los tramos $\mathcal{D}$, y
$h^{sr}_{k,g}(\cdot)$ una función escalonada en tramos quinquenales de edad
(anidados en los macro-tramos, como controles de nivel). La interacción
tipo$\times$duración es el objeto central —capturar cuánta memoria de
duración sobrevive dentro de tipo— y la interacción
duración$\times$macro-tramo permite que esa memoria cambie a lo largo del
ciclo de vida: la restricción de separabilidad edad--duración falla de
forma económicamente relevante justo en el tipo III a edades altas
(hasta 10 puntos porcentuales en la probabilidad de seguir en laguna a
cinco años), donde las lagunas se vuelven cuasi-absorbentes.

\begin{assumption}[Exogeneidad del proceso de participación]
\label{as:exo}
$\{s_{it}\}$ es exógeno respecto de las decisiones de inversión del fondo y
de la política de carteras de referencia.
\end{assumption}

El Supuesto~\ref{as:exo} es estándar en la literatura de ciclo de vida
(participación laboral no responde a la rentabilidad del ahorro forzoso) y
es, si acaso, conservador para nuestro propósito de política.

\subsection{Representación markoviana en el estado ampliado}

Definiendo el estado ampliado $x_{it} = (s_{it}, d_{it}) \in \mathcal{X}$ con
$|\mathcal{X}| = 2\bar D$ (o, de forma equivalente y más compacta para la
programación dinámica, sobre los tramos de duración $\mathcal{D}$ con un
contador de meses dentro del tramo), el proceso es Markov de primer orden
condicional a $(k, a, g)$, con kernel
\begin{equation}
P_{k,g,a}\bigl((s,d) \to (s', d')\bigr) =
\begin{cases}
\lambda^{s\,(1-s)}_{k,g}(d,a) & \text{si } s' = 1-s,\ d' = 1,\\[2pt]
1 - \lambda^{s\,(1-s)}_{k,g}(d,a) & \text{si } s' = s,\ d' = \min\{d+1,\bar D\},\\[2pt]
0 & \text{en otro caso.}
\end{cases}
\label{eq:kernel}
\end{equation}
Esta representación es la que ingresa al modelo de programación dinámica: el
costo de la memoria de duración es lineal en $\bar D$, no exponencial.

\subsection{Condiciones iniciales}

A la edad de entrada $a_0$ (afiliación), el individuo extrae
$x_{i,0} = (s_0, 1)$ con $\Prob(s_0 = 1 \mid k, g)$ estimada de la
distribución empírica del primer año post-afiliación. Las proporciones de
tipos por sexo $\pi_{k|g}$ se estiman en la clasificación
(Sección~\ref{sec:estimacion}).

\section{El proceso de remuneraciones}

\subsection{Especificación}

Toda remuneración se expresa en UF (unidades de fomento) del mes
correspondiente, de modo que el proceso está en términos reales por
construcción. Condicional a cotizar ($s_{it}=1$), la remuneración imponible
observada es
\begin{equation}
w^{obs}_{it} = \min\{w_{it},\ \bar w_t\}, \qquad
\log w_{it} = m_{k,g}(a_{it}) + \zeta_{i} + z_{it} + \varepsilon_{it},
\label{eq:wage}
\end{equation}
donde $\bar w_t$ es el tope imponible vigente (censura observable, con
indicador administrativo en los datos), $m_{k,g}(\cdot)$ es el perfil
edad--ingreso por tipo y sexo, $\zeta_{i}$ un efecto fijo individual,
$z_{it}$ el componente persistente y
$\varepsilon_{it}$ ruido transitorio (que incluye error de medición y
fracción de mes trabajada):
\begin{equation}
z_{it} = \rho_{k,g}\, z_{i,t-1} + \eta_{it}, \qquad
\eta_{it} \sim iid\,(0, \sigma^2_{\eta,k,g}), \qquad
\varepsilon_{it} \sim iid\,(0, \sigma^2_{\varepsilon,k,g}),
\end{equation}
con todos los parámetros del proceso específicos al tipo y al sexo.

\begin{assumption}[Identificación edad--período--cohorte]
\label{as:apc}
La tendencia común del ingreso real se atribuye a efectos de período
(crecimiento agregado de la productividad) y no a efectos de cohorte. El
perfil de edad $m_{k,g}(\cdot)$ se identifica de la variación intra-individuo
neta de efectos de tiempo; su nivel se ancla a los años recientes
(2015--2023); y el crecimiento salarial real futuro se introduce en el modelo
de ciclo de vida como un parámetro explícito $g$, sujeto a análisis de
sensibilidad (rango de referencia 1,25\%--2\% anual, consistente con las
proyecciones oficiales).
\end{assumption}

El Supuesto~\ref{as:apc} responde a la fuerte no-estacionariedad de la
muestra: la remuneración real mediana de los cotizantes cae $\sim$35\% entre
1981 y 1987 y luego casi se triplica ($\sim$2,4\% real anual hasta 2023),
mientras la composición etaria del panel también se desplaza. Sin una
normalización explícita, la identificación de $m_{k,g}(a)$ es imposible
(edad $=$ cohorte $+$ tiempo; \citealp{deaton1994}) y el crecimiento agregado
se confundiría con retornos a la edad. La atribución a período y no a cohorte
es el supuesto estándar en la literatura de ciclo de vida
\citep{cocco2005} y es discutible;
se reportará la sensibilidad del perfil a la normalización alternativa.

\begin{assumption}[Neutralidad salarial de las lagunas]
\label{as:noscar}
El componente persistente $z_{it}$ evoluciona en tiempo calendario con la
misma ley de movimiento durante episodios de cotización y de laguna; la
remuneración es simplemente no observada cuando $s_{it}=0$.
\end{assumption}

El Supuesto~\ref{as:noscar} está respaldado por el hecho estilizado 4: el
salario mediano de reingreso, deflactado y neto del crecimiento contrafactual,
no exhibe penalización por duración de la laguna (las estimaciones puntuales
son levemente positivas, consistentes con selección en el reingreso). Esto
contrasta con la práctica de imponer depreciación de capital humano durante
las interrupciones y simplifica el estado del modelo: no se requiere
arrastrar la duración acumulada de lagunas como determinante salarial.

\subsection{Cotización}

La cotización mensual al fondo es
\begin{equation}
c_{it} = \tau_t \cdot \min\{w_{it}, \bar w_t\} \cdot s_{it},
\end{equation}
con $\tau_t$ la tasa legal vigente (10\% histórico; trayectoria creciente
según la Ley 21.735 para el aporte con destino a cuenta individual). Tanto
$\tau_t$ como $\bar w_t$ son parámetros institucionales deterministas, no
objetos estocásticos.

\section{Estimación}
\label{sec:estimacion}

Toda la estimación se realiza separadamente por sexo. La estrategia tiene
cuatro pasos:

\paragraph{Paso 1: tipos por mixture de verosimilitud completa.} Los tipos
se estiman por EM \citep{heckman1984}: el tipo es la única variable latente,
de modo que, condicional a $k$, la verosimilitud de la secuencia mensual del
individuo es un producto de Bernoullis sobre sus transiciones, y el
estadístico suficiente individual son sus conteos por celda
estado$\times$duración$\times$edad. El paso E entrega probabilidades
posteriores de tipo; el paso M re-estima los logits \eqref{eq:logit} con
pesos posteriores, junto con $\pi_{k|g}$ y las condiciones iniciales. $K$ se
elige comparando $K=2,\dots,5$ por sexo bajo la especificación
\eqref{eq:logit}: el BIC decrece monótonamente (sobre-selección esperable
con cientos de observaciones por individuo), el ICL —que penaliza
clasificaciones difusas— alcanza su mínimo en $K=3$ para hombres, y para
mujeres queda esencialmente empatado entre $K=2$ y $K=4$ (11 puntos sobre
$1{,}06$ millones) con $K=3$ cercano; el $K=4$ femenino solo agrega una
clase degenerada de 2,5\%. Se adopta $K=3$ como especificación base en
ambos sexos ($K=2$ y $K=4$ como robustez). Las proporciones estimadas son
desiguales ($\hat\pi_{k|g}$: tipo I 0,48 y 0,53; tipo II 0,18 y 0,14;
tipo III 0,34 y 0,34, para hombres y mujeres respectivamente). La
clasificación dura por
terciles de densidad —usada como inicialización y robustez— coincide con
el tipo modal EM en los polos pero no en el centro: el EM agrupa por
\emph{dinámica} de los hazards y no por densidad bruta, y reasigna una
fracción sustancial del tercil medio de densidad al tipo I. Por lo
mismo, el tipo III no equivale a ``casi no cotiza'': sus densidades
posteriores medias son 0,33 (hombres) y 0,26 (mujeres). Menos de la mitad
de los individuos exhibe posterior modal superior a 0,9, lo que desaconseja
la clasificación dura como especificación principal.

Un punto metodológico que corresponde declarar: la verosimilitud del
mixture es multimodal. En mujeres, una búsqueda multi-arranque encuentra un
óptimo local alternativo con verosimilitud levemente superior
($+64$ puntos sobre $523$ mil, un 0,01\%) que reorganiza la tipología —
concentra el 62\% de las mujeres en una sola clase que mezcla dinámicas de
adhesión estable con historias de densidad muy baja. Entre óptimos locales
prácticamente empatados en verosimilitud, la selección se resolvió por
momentos \emph{no utilizados} en la estimación (Sección de validación): la
solución adoptada domina a la alternativa en todos los momentos de
validación —en particular, la alternativa deteriora la cola baja de la
distribución de densidades (0,129 contra 0,146, con 0,155 en los datos)— y
la alternativa se reporta como robustez.

\paragraph{Paso 2: hazards.} Logits \eqref{eq:logit} por tipo, sexo,
estado de origen y macro-tramo de edad, con dummies de duración sobre los
tramos $\mathcal{D}$ y tramos quinquenales de edad anidados; la estimación
es el paso M del EM (MLE exacto sobre celdas agregadas ponderadas por
posteriores, vía el estadístico suficiente). La especificación se
seleccionó comparando alternativas anidadas (truncada en 24; duraciones
extendidas aditiva; duraciones extendidas con macro-tramos) no por
contrastes de razón de verosimilitud —que con $7{,}5$ millones de
observaciones rechazan cualquier restricción— sino por su desempeño en los
objetos que el modelo de ciclo de vida consume: curvas de supervivencia en
laguna por edad de inicio y cotizaciones esperadas a 60 meses condicionales
a (edad, estado, duración). La versión truncada en 24 sobreestimaba las
cotizaciones esperadas de los afiliados de 55--59 años en laguna en 5 a 11
meses (sobre niveles de 7 a 16); la especificación adoptada reduce ese
error a 0,5--3 meses. Robustez temporal obligatoria: re-estimación por
subperíodos (1990--2005 y 2008--2023); ante inestabilidad estructural se
privilegia el período reciente, que es el relevante para la proyección.

\paragraph{Paso 3: remuneraciones.} Con la submuestra $s_{it}=1$,
excluyendo observaciones en el tope (flag administrativo), los meses cuya
remuneración incluye subsidios de incapacidad (la base de la licencia no es
salario de mercado) y, como robustez, los meses con registros repetidos del
mismo pagador (pagos retroactivos que inflan el mes de entero): (i)
$m_{k,g}(a)$ y $\zeta_{i}$ por regresión con efectos fijos individuales sobre
splines de edad \emph{y efectos de año}, por tipo y sexo, bajo la
normalización del Supuesto~\ref{as:apc}, con nivel anclado a 2015--2023;
(ii) $\rho_{k,g}$, $\sigma^2_{\eta,k,g}$ y $\sigma^2_{\varepsilon,k,g}$ por GMM sobre
autocovarianzas de los residuos en frecuencia anual (promedio de meses
cotizados del año), siguiendo la práctica estándar para separar componente
persistente y transitorio. La censura por tope se trata por máxima
verosimilitud tobit como robustez.

\paragraph{Paso 4: validación por momentos no utilizados.} El proceso
estimado se simula (mensualmente, desde la edad de afiliación) y se contrasta
con momentos \emph{no} usados en la estimación:
distribución completa de densidades individuales (la forma de U),
distribución de duración de lagunas completadas, densidad por edad y sexo,
número de episodios de laguna por individuo a los 60 años, autocorrelación
de la densidad anual, y —dado su papel en la selección de especificación—
los momentos \emph{prospectivos}: meses cotizados efectivos en los 60 meses
siguientes, condicionales a la celda (edad, estado, duración) de partida. La agregación mensual$\to$anual para el modelo de
ciclo de vida se efectúa sobre estas mismas simulaciones, eligiendo la
representación anual más parsimoniosa que preserve los momentos relevantes
para la acumulación de saldo.

\section{Resultados de la estimación}
\label{sec:resultados}

Con $K=3$, los hazards (2 sexos $\times$ 3 tipos $\times$ 2 estados
$\times$ 3 macro-tramos de edad) provienen del paso M del EM en el óptimo:
celdas agregadas duración$\times$edad ponderadas por las probabilidades
posteriores de tipo. La Figura~\ref{fig:hazfit} muestra el ajuste para la
salida de laguna: el modelo reproduce los niveles —que difieren por un
factor de 3 entre tipos—, la caída con la duración —que continúa hasta las
lagunas de 120+ meses— y el máximo local en los 9--11 meses. El ordenamiento de niveles se invierte en los hazards de
término de cotización, confirmando que los tipos capturan dinámicas espejo y
no solo niveles. La robustez por subperíodo indica estabilidad razonable:
los hazards 2008--2023 son en promedio 6\% más altos (en logs, $+0{,}06$)
que los de 1990--2005, con la reentrada desde laguna $\sim$11\% más rápida
en el período reciente.

\begin{figure}[htbp]
\centering
\includegraphics[width=\textwidth]{../output/calibration/fig8_ajuste_hazard_laguna.png}
\caption{Hazard mensual de salida de laguna por duración: celdas empíricas
(puntos) y ajuste del logit (líneas), por tipo de adhesión y sexo, promediado
sobre tramos de edad con pesos muestrales.}
\label{fig:hazfit}
\end{figure}

\subsection{La firma dinámica de los tipos}
\label{sec:firmas}

La Figura~\ref{fig:firmas} traduce los hazards a su objeto interpretable:
la probabilidad de que un episodio de empleo o de laguna, iniciado a los 30
años, sobreviva más de $d$ meses. El primer orden de la estructura es una
dicotomía de persistencia alta/baja en cada estado, que ubica a los tipos
en tres esquinas de un espacio dos por dos. El tipo I (estable) es
persistencia alta en empleo y baja en laguna: empleos con mediana de 14
meses (hombres) y 10 (mujeres), más de un quinto sobre los cinco años, y
lagunas con mediana de 3 meses. El tipo II (intermitente) es persistencia
baja en ambos estados, con colas delgadas: episodios de 3--4 meses de
mediana y $\Prob(\text{episodio} > 5\text{ años})$ inferior al 4\% en ambos
estados, lo que —por promediación sobre muchos episodios— comprime sus
densidades de vida hacia el centro de la distribución. El tipo III combina
empleos de mediana corta (3 meses) con las lagunas más persistentes de la
población (mediana de 10 meses en hombres y 22 en mujeres;
$\Prob(>5\text{ años})$ de 0,14 y 0,31).

Esta taxonomía no está impuesta entre sexos. La estimación es completamente
independiente para hombres y mujeres —verosimilitudes, hazards y
proporciones separados, sin ningún parámetro común— y nada obliga a que los
mixtures de ambos sexos encuentren perfiles comparables: condicional a
$K=3$, podrían haber emergido seis firmas sin relación entre sí. Emergen,
en cambio, las mismas tres firmas cualitativas en ambos sexos, con
proporciones además muy próximas —en el caso del tipo III, casi idénticas
(0,337 y 0,340)—, lo que sugiere que la tipología captura una estructura
real del empleo formal chileno y no un artefacto muestral o del algoritmo.
Las diferencias entre sexos se cargan a la \emph{severidad} de la firma del
tipo III (lagunas más largas y más tempranas en mujeres) y a la repartición
entre los tipos I y II, no a la existencia de los tipos. Como se discute
más abajo, esta simetría es propia del nivel $K=3$: con $K=4$ las
estructuras de ambos sexos divergen.

El segundo orden está en las colas, y es lo que hace al tipo III resistente
a cualquier etiqueta de nivel. Aunque su empleo mediano es tan corto como
el del tipo II, su cola de empleo es cuatro a seis veces más gruesa
($\Prob(>5\text{ años}) \approx 0{,}09$--$0{,}10$), y por sesgo de longitud
esos empleos largos concentran el 72--75\% de sus meses efectivamente
cotizados —la misma proporción que en el tipo I—. El resultado es
polarización: el tipo III genera tanto vidas de densidad casi nula como
vidas casi completas. Es además el único tipo cuya firma cambia
sustancialmente con la edad, con asimetría por sexo: en hombres,
$\Prob(\text{laguna nueva} > 5\text{ años})$ sube de 0,14 a los 30 a 0,25
a los 50; en mujeres, la laguna más persistente es la iniciada a los 30
—consistente con interrupciones en edad de crianza que no se revierten—.

La cuarta esquina del espacio —persistencia alta en ambos estados— no está
vacía sino escondida. Re-estimando el mixture con $K=4$, en hombres esa
esquina emerge como una clase del 27\% (cola de empleo al nivel del tipo I;
lagunas aún más persistentes que las del tipo III con $K=3$) y la esquina
de empleo corto con laguna larga casi se vacía: el tipo III masculino es en
buena medida un compuesto de carreras largas que terminan en salidas
prolongadas o definitivas. En mujeres, en cambio, $K=4$ no produce esa
esquina (solo una clase degenerada de 2,5\%): el tipo III femenino sí es
el espejo del tipo I. No se adopta $K=4$ como especificación base por tres
razones: el ICL prefiere $K=3$ también en hombres, la estructura de cuatro
tipos no existe en ambos sexos, y la comparabilidad entre sexos es central
para el uso de política del proceso. El costo de esa decisión
—heterogeneidad residual dentro del tipo III— queda declarado y
cuantificado (momentos prospectivos y colas de la validación), y la
especificación masculina con $K=4$ se reporta como robustez.

\begin{figure}[htbp]
\centering
\includegraphics[width=\textwidth]{../output/calibration/fig21_firmas_tipos.png}
\caption{La firma de cada tipo: probabilidad de que un episodio de empleo
(izquierda) o de laguna (derecha) sobreviva más de $d$ meses, para
episodios iniciados a los 30 años (líneas llenas) y a los 50 (punteadas),
según los hazards estimados. Escala horizontal logarítmica.}
\label{fig:firmas}
\end{figure}

\subsection{El proceso de remuneraciones estimado}

La Figura~\ref{fig:perfiles} presenta los perfiles edad--ingreso anclados a
2015--2023. Los efectos de período absorbieron un crecimiento común de
2,7--4,0\% anual —superior al índice agregado de remuneraciones reales, lo
que refleja composición del empleo formal— y que en el modelo de ciclo de
vida entrará como el parámetro $g$. Las varianzas anuales (netas del efecto
fijo individual) arrojan $\rho_{k,g} \in [0{,}50,\ 0{,}68]$ y
$\sigma^2_{\eta} \in [0{,}05,\ 0{,}09]$ en los seis grupos. La
especificación AR(1) es consistente con los datos al horizonte observable:
las razones de autocovarianzas sucesivas ($\hat c_2/\hat c_1$ y
$\hat c_3/\hat c_2$) son aproximadamente iguales entre sí y al $\hat\rho$
implícito en los seis grupos, como exige el decaimiento geométrico. La
única estimación restringida es la del tipo II femenino, donde $\hat\rho$
queda en la cota inferior impuesta de 0,5 (las razones de autocovarianzas
sugieren $\rho \approx 0{,}45$; es el grupo con menos persona--año y mayor
error de medición por promediar pocos meses cotizados). Un dividendo de la
clasificación EM: bajo terciles, el grupo femenino del tipo III era inestimable
(6.242 persona--año y selección severa); con tipos EM reúne 16.380
persona--año y sus parámetros son regulares ($\rho = 0{,}65$), eliminando la
necesidad de agrupación ad hoc.

\begin{figure}[htbp]
\centering
\includegraphics[width=\textwidth]{../output/calibration/fig10_perfiles_edad_ingreso.png}
\caption{Perfiles edad--ingreso $\exp(\hat m_{k,g}(a))$ por tipo y sexo, en
UF mensuales al nivel 2015--2023 (escala log). El tramo femenino sobre los 60
años refleja selección post edad legal de retiro y no se utiliza en el
modelo.}
\label{fig:perfiles}
\end{figure}

El ordenamiento no monótono de la Figura~\ref{fig:perfiles} —el tipo III
exhibe remuneraciones observadas iguales o superiores a las del tipo I, y
el tipo II muy inferiores a ambos— es informativo y esperable: la
remuneración solo se observa \emph{condicional a cotizar}. Para el tipo
II, la alta rotación y la baja remuneración son el mismo fenómeno
(empleos formales cortos y mal pagados). Para el tipo III, los escasos
meses cotizados son una muestra seleccionada de sus mejores estados (lógica
de salario de reserva), selección que se intensifica con la edad. Esta
selección no sesga el uso del proceso en el modelo de ciclo de vida —la
acumulación de saldo requiere precisamente la remuneración condicional a
cotizar— pero sería central si la participación se endogenizara.

\subsection{Validación con momentos no utilizados}

El proceso completo se simuló replicando la estructura exacta del panel
(ventana observada, tipo y sexo de cada individuo; semilla fija). El Cuadro
\ref{tab:valid} y la Figura~\ref{fig:validU} comparan momentos que la
estimación nunca utilizó. La forma de U de las densidades individuales
—el hecho estilizado central— emerge de la dinámica estimada, y la
distribución de duración de lagunas queda esencialmente clavada (media
11,5 vs.\ 11,3 meses; p90 de 28 vs.\ 27). Los desajustes residuales son
pequeños, simétricos y de signo declarado: ambas colas de la U quedan
levemente subpredichas (cola baja 0,146 vs.\ 0,155; cola alta 0,148 vs.\
0,165). La subpredicción de la cola alta —que reduce los pensionados
autofinanciados simulados y por tanto \emph{no} es conservadora para la
aplicación de política— era de 2,6 puntos porcentuales bajo la
especificación truncada en $\bar D = 24$ y las duraciones extendidas la
reducen a 1,7; se tendrá presente al interpretar los resultados del modelo
de ciclo de vida.

Una precisión sobre el estatus de estos momentos. Todos son externos a la
verosimilitud, pero no todos son ajenos a toda decisión: la distribución de
densidades individuales participó en la selección entre óptimos locales del
EM (Paso 1), y las curvas de supervivencia y los momentos prospectivos
guiaron la elección de la especificación de duración, de modo que su papel
aquí es de consistencia y no de validación ciega. Permanecen fuera de toda
decisión de estimación o especificación —y constituyen por tanto la
validación estricta— la densidad por edad y sexo, el número de episodios de
laguna por individuo, la mediana de duración de lagunas y la
autocorrelación de la densidad anual; esta última es el test más exigente
de la dinámica de mediano plazo y el modelo la reproduce con leve
sub-persistencia (0,825 vs.\ 0,842 al primer rezago), consistente con la
heterogeneidad intra-tipo residual discutida en las limitaciones.

\begin{table}[htbp]
\centering
\begin{tabular}{lcc}
\toprule
 & Datos & Modelo \\
\midrule
Densidad media                          & 0,537 & 0,537 \\
Densidad individual: p10                & 0,040 & 0,055 \\
\quad p50                               & 0,563 & 0,561 \\
\quad p90                               & 0,953 & 0,937 \\
Share densidad $<$ 0,10                 & 0,155 & 0,146 \\
Share densidad $>$ 0,90                 & 0,165 & 0,148 \\
Duración de lagunas: mediana (meses)    & 3     & 3     \\
\quad media                             & 11,3  & 11,5  \\
\quad p90                               & 27    & 28    \\
Episodios de laguna por individuo       & 7,6   & 7,4   \\
Autocorr.\ densidad anual: rezago 1     & 0,842 & 0,825 \\
\quad rezago 2                          & 0,733 & 0,700 \\
\bottomrule
\end{tabular}
\vspace{4pt}
\caption{Momentos no utilizados en la estimación: datos vs.\ modelo}
\label{tab:valid}
\end{table}

\begin{figure}[htbp]
\centering
\includegraphics[width=0.75\textwidth]{../output/calibration/fig12_validacion_U.png}
\caption{Distribución de densidades individuales de cotización: datos
(histograma) vs.\ proceso estimado simulado sobre las mismas ventanas
(línea). Momento no utilizado en la estimación.}
\label{fig:validU}
\end{figure}

\section{Limitaciones y decisiones abiertas}

\begin{enumerate}
\item \textbf{Muestra.} La HPA proviene de la muestra teórica de la EPS;
sobrerrepresenta cohortes antiguas y no incluye trabajadores nunca afiliados.
El ``chileno mediano'' del producto 2 se definirá sobre la población
\emph{afiliada}, que es la relevante para el glide path.
\item \textbf{Estacionalidad.} Ambos hazards exhiben un máximo local en
duraciones de 9--11 meses (contratos anuales, trabajo estacional). Las
dummies de duración lo absorben; no se modela estructuralmente.
\item \textbf{Identificación duración vs.\ heterogeneidad residual.} Dentro
de tipo puede persistir heterogeneidad; la lectura causal de $f^{sr}_k$ no es
necesaria para nuestro propósito (simulación de trayectorias), pero sí lo es
para interpretaciones estructurales, y así se declarará. La huella
cuantitativa de esa heterogeneidad residual es visible en los momentos
prospectivos: \emph{toda} ley de primer orden en $(s,d)$ —incluidas las
celdas saturadas— sobrepredice en 1 a 3 meses las cotizaciones efectivas a
60 meses de quienes se observan en lagunas de duración media a edades
jóvenes (selección dinámica: quien lleva dos años en laguna a los 27 tiene
peor futuro del que su celda indica). La robustez con $K=4$ acota cuánto de
esto absorbería un tipo adicional.
\item \textbf{Edad de término.} El panel se censura en la solicitud de
pensión o los 65 años; la extensión post-65 (relevante para el Fondo de
Consolidación) requerirá supuestos adicionales.
\item \textbf{Multimodalidad del EM.} La selección entre óptimos locales
del mixture se resolvió por validación con momentos no utilizados (véase el
Paso 1); la tipología femenina es sensible al óptimo elegido y la
alternativa se reporta como robustez. Es una limitación conocida de los
modelos de tipos discretos à la Heckman--Singer, aquí explícita en lugar de
oculta tras una única inicialización.
\item \textbf{Pendiente.} La frecuencia del modelo de ciclo de vida
(mensual vs.\ anual con proceso agregado); $K$ y la estructura de duración
quedaron fijados ($K=3$ por ICL y comparabilidad entre sexos; tramos de
duración extendidos hasta 120+ meses con perfiles por macro-tramo de edad).
\end{enumerate}

\bibliographystyle{chicago}
\bibliography{referencias}

\end{document}
